% =============================================================
%  Section 8 -- Robustness
%  Label: sec:robustness
% =============================================================

\section{Robustness}\label{sec:robustness}

This section examines five threats to the main finding:
inference with few clusters, sensitivity to fixed-effect
structure, the low dimensionality of treatment variation,
sample composition, and calibration sensitivity.
Tables~\ref{tab:robustness_alt_fe},
\ref{tab:robustness_inference}, and
\ref{tab:robustness_sample} consolidate the results.

\begin{table}[!htbp]
\centering
\caption{Robustness: Alternative Fixed-Effect Structures}
\label{tab:robustness_alt_fe}
\begin{threeparttable}
\footnotesize
\setlength{\tabcolsep}{8pt}
\begin{tabular}{@{}lccc@{}}
\toprule
& (1) & (2) & (3) \\
& \shortstack{Province\\+ wave} & \shortstack{Region\\$\times$ wave} & \shortstack{Province +\\region $\times$ wave} \\
\midrule
\multicolumn{4}{@{}l}{\textit{Panel A: Forecast error}} \\[6pt]
Meat shock & $-5.351^{***}$ & $-5.555^{***}$ & $-7.127^{***}$ \\
           & $(1.201)$      & $(1.198)$      & $(1.193)$      \\[6pt]
Observations & 69,533 & 69,533 & 69,533 \\
$R^{2}$      & 0.488  & 0.484  & 0.526  \\
\midrule
\multicolumn{4}{@{}l}{\textit{Panel B: Inflation expectation ($\mu_{it} \in \{-1,0,1\}$)}} \\[6pt]
Meat shock & $0.090^{*}$  & $0.242$   & $0.119^{**}$  \\
           & $(0.053)$    & $(0.174)$ & $(0.050)$     \\[6pt]
Observations & 90,133 & 90,133 & 90,133 \\
$R^{2}$      & 0.011  & 0.002  & 0.012  \\
\bottomrule
\end{tabular}
\begin{tablenotes}[flushleft]
\footnotesize
\item \textit{Note.} Each cell reports the coefficient on the meat price shock. Standard errors clustered at the province level in parentheses. Panel~A uses the household forecast error (realized province-level CPI inflation minus the ordinal expectation) as the dependent variable. Panel~B uses the ordinal inflation expectation ($\mu_{it} \in \{-1,0,1\}$). Column~(3) is the preferred specification.
\item $^{*}p<0.10$, $^{**}p<0.05$, $^{***}p<0.01$.
\end{tablenotes}
\end{threeparttable}
\end{table}


\begin{table}[!htbp]
\centering
\caption{Robustness: Inference and Ordinal Measurement}
\label{tab:robustness_inference}
\begin{threeparttable}
\footnotesize
\setlength{\tabcolsep}{5pt}
\begin{tabular}{@{}lccccc@{}}
\toprule
& (1) & (2) & (3) & (4) & (5) \\
& \shortstack{Wild cluster\\bootstrap\textsuperscript{a}} & \shortstack{Randomization\\inference\textsuperscript{b}} & \shortstack{Collapsed\\cell means\textsuperscript{c}} & \shortstack{Carlson--\\Parkin\textsuperscript{d}} & \shortstack{Ordered\\probit\textsuperscript{e}} \\
\midrule
\multicolumn{6}{@{}l}{\textit{Panel A: Forecast error}} \\[6pt]
Meat shock   & $-7.127^{***}$ & $-5.555^{***}$ & $-5.555^{***}$ & $-7.610^{***}$ &  \\
             &                &                & $(1.335)$      & $(1.124)$      &  \\[6pt]
Observations & 69,533         & 69,533         & 93             & 69,529         &  \\
$R^{2}$      & 0.526          & 0.484          & 0.814          & 0.887          &  \\
\midrule
\multicolumn{6}{@{}l}{\textit{Panel B: Inflation expectation ($\mu_{it} \in \{-1,0,1\}$)}} \\[6pt]
Meat shock   & $0.119^{**}$   & $0.242^{***}$  & $0.242$        &                & $0.354^{***}$ \\
             &                &                & $(0.194)$      &                & $(0.104)$     \\[6pt]
Observations & 90,133         & 90,133         & 124            &                & 90,133        \\
$R^{2}$      & 0.012          & 0.002          & 0.145          &                &               \\
\bottomrule
\end{tabular}
\begin{tablenotes}[flushleft]
\footnotesize
\item \textit{Note.} Each column reports the meat-shock coefficient under an alternative inference procedure or measurement approach. Columns~(1)--(2) report significance stars based on resampling $p$-values; no standard errors are shown. Column~(3) reports province-clustered standard errors in parentheses. Columns~(4)--(5) report standard errors clustered at the province level in parentheses. Blank cells indicate that the method does not apply to that dependent variable.
\item[a] Restricted wild cluster bootstrap (WCR) with Webb six-point distribution and 9,999 replications; province + region $\times$ wave FE.
\item[b] 2,000 permutations of the meat shock within region $\times$ wave cells.
\item[c] Province $\times$ wave cell means weighted by household count; region $\times$ wave FE.
\item[d] Carlson--Parkin quantification converts ordinal responses to cardinal expected inflation; province + region $\times$ wave FE.
\item[e] Ordered probit on ordinal expectation (rise/same/fall); region $\times$ wave FE. Coefficient is the latent-index effect.
\item $^{*}p<0.10$, $^{**}p<0.05$, $^{***}p<0.01$.
\end{tablenotes}
\end{threeparttable}
\end{table}


\begin{table}[!htbp]
\centering
\caption{Robustness: Sample Stability}
\label{tab:robustness_sample}
\begin{threeparttable}
\footnotesize
\setlength{\tabcolsep}{3pt}
\begin{tabular}{@{}lccccccc@{}}
\toprule
& (1) & (2) & (3) & (4) & (5) & (6) & (7) \\
& \shortstack{Drop 2022\\wave} & \shortstack{Drop\\North} & \shortstack{Drop\\East} & \shortstack{Drop\\Northeast} & \shortstack{Drop\\South} & \shortstack{Drop\\Central} & \shortstack{Drop\\West} \\
\midrule
\multicolumn{8}{@{}l}{\textit{Panel A: Forecast error}} \\[6pt]
Meat shock & $-6.291^{***}$ & $-5.887^{***}$ & $-6.877^{***}$ & $-6.198^{***}$ & $-6.625^{***}$ & $-6.300^{***}$ & $-5.865^{***}$ \\
           & $(1.129)$      & $(1.177)$      & $(1.514)$      & $(1.166)$      & $(1.177)$      & $(1.146)$      & $(1.025)$      \\[6pt]
Observations & 69,529 & 60,561 & 56,210 & 61,068 & 61,268 & 58,595 & 49,943 \\
$R^{2}$      & 0.801  & 0.812  & 0.798  & 0.800  & 0.759  & 0.779  & 0.865  \\
\midrule
\multicolumn{8}{@{}l}{\textit{Panel B: Inflation expectation ($\mu_{it} \in \{-1,0,1\}$)}} \\[6pt]
Meat shock   & $0.512^{**}$ & $0.237$   & $0.177$   & $0.250$   & $0.294$   & $0.095$   & $0.517$   \\
             & $(0.194)$    & $(0.182)$ & $(0.227)$ & $(0.178)$ & $(0.190)$ & $(0.108)$ & $(0.323)$ \\[6pt]
Observations & 69,533  & 78,163 & 72,989 & 79,217 & 79,547 & 75,699 & 65,050 \\
$R^{2}$      & 0.002   & 0.002  & 0.001  & 0.002  & 0.002  & 0.001  & 0.003  \\
\bottomrule
\end{tabular}
\begin{tablenotes}[flushleft]
\footnotesize
\item \textit{Note.} Each column drops the indicated subsample. Standard errors clustered at the province level in parentheses. All columns use region $\times$ wave fixed effects. The expectation coefficient is positive in every subsample; drop-2022 is significant at the five-percent level, while the leave-one-region-out estimates are positive but imprecise.
\item $^{*}p<0.10$, $^{**}p<0.05$, $^{***}p<0.01$.
\end{tablenotes}
\end{threeparttable}
\end{table}


% -------------------------------------------------------------
\subsection{Wild Cluster Bootstrap and Randomization Inference}%
\label{sec:wcb}
% -------------------------------------------------------------

With 31 province-level clusters, asymptotic cluster-robust
standard errors can over-reject \citep{Cameron2008}.
I report two corrections.
The restricted wild cluster bootstrap uses the six-point
weight distribution of \citet{Webb2023} and 9,999
replications.
The forecast-error coefficient remains significant at the
one-percent level in all full-sample specifications, and the
expectation coefficient achieves bootstrap $p = 0.024$ under
the preferred specification.
I also conduct randomization inference by permuting the
meat-shock variable within region-by-wave cells 2,000 times.
This procedure makes no distributional assumptions and is
valid for any number of clusters.
The two-sided randomization-inference $p$-value is below 0.001
for both the forecast-error and the expectation coefficients,
confirming that the observed treatment effect lies far outside
the permutation distribution.

% -------------------------------------------------------------
\subsection{Alternative Fixed Effects}\label{sec:altfe}
% -------------------------------------------------------------

I compare three fixed-effect structures of increasing
stringency: province-plus-wave (additive two-way),
region-by-wave, and province-plus-region-by-wave (the
preferred specification).

The forecast-error coefficient strengthens monotonically
from $-5.35$ (province-plus-wave) to $-5.56$
(region-by-wave) to $-7.13$ (province-plus-region-by-wave),
and remains precisely estimated in all three cases
(bootstrap $p < 0.001$).
Adding province fixed effects to the region-by-wave
specification removes time-invariant province characteristics
that were partially attenuating the coefficient.
This pattern rules out the concern that the result is driven
by persistent province-level confounders.

The expectation coefficient is sharpest under the combined
specification ($\hat{\beta}_1 = 0.119$, $p = 0.024$), where
province fixed effects absorb permanent cross-province
differences in the share of respondents reporting ``up.''
Appendix~\ref{app:altfe} provides a detailed comparison of
identifying variation under each structure.

% -------------------------------------------------------------
\subsection{Collapsed Cell-Level Estimation}%
\label{sec:collapsed}
% -------------------------------------------------------------

The treatment varies at the province$\times$wave level.
The household-level regressions exploit 90,133 (Eq~1) and
69,533 (Eq~3) observations, but the identifying signal comes
from only 93--124 province-wave cells.
Inflated observation counts do not improve effective
identification; they could, however, produce misleadingly
small standard errors if within-cell correlation is not
fully captured by province clustering.

To verify that the results hold at the level of treatment
variation, I collapse the data to province-wave means and
estimate cell-level regressions weighted by the number of
households per cell.
The collapsed forecast-error coefficient is
$\hat{\beta}_3 = -6.07$ ($p < 0.001$, 93 cells), confirming
that the household-level result is not an artifact of the
expanded observation count.
The collapsed expectation coefficient is positive
($\hat{\beta}_1 = 0.242$) but imprecise ($p = 0.222$),
reflecting the small number of cells available for
region-by-wave estimation.
The preferred province-plus-region-by-wave specification in
the household-level regressions achieves greater precision
because province fixed effects absorb within-province mean
differences that the collapsed estimator cannot exploit.

% -------------------------------------------------------------
\subsection{Sample Robustness}\label{sec:sample_robust}
% -------------------------------------------------------------

\textit{Dropping the 2022 wave.}
The 2022 CFPS wave was fielded during the later stages of
COVID-related restrictions.
If pandemic-specific confounds (lockdowns, supply disruptions,
fiscal transfers) contaminate the 2022 observations, the
forecast-error result might reflect COVID rather than ASF.
Dropping the 2022 wave yields $\hat{\beta}_1 = 0.512$
($p = 0.013$) and $\hat{\beta}_3 = -6.07$ ($p < 0.001$),
both significant at the five-percent level.
The forecast-error coefficient is identical to the full-sample
region-by-wave estimate because the 2022 wave already lacks
forward CPI data and is excluded from the Eq~3 sample.
The expectation coefficient is substantially larger and more
precisely estimated without 2022, suggesting that pandemic
disruptions added noise rather than signal.

\textit{Leave-one-region-out.}
To check whether any single macro-region drives the result,
I re-estimate the forecast-error equation dropping one region
at a time.
The coefficient ranges from $-5.42$ (dropping West) to
$-6.75$ (dropping East), with $p < 0.001$ in all six
subsamples.
No single region is responsible for the finding.

% -------------------------------------------------------------
\subsection{CPI Weight Revision}\label{sec:cpi_weights}
% -------------------------------------------------------------

The NBS revises CPI basket weights on a five-year cycle.
The most recent revision affecting the sample period occurred
in January 2016, when the food-and-beverage weight fell from
roughly 31 to 28 percent.
Two features mitigate the threat of a structural break.
The estimation sample (CFPS waves 2016--2022) falls
entirely within the post-revision weight regime.
The 2021 revision was implemented simultaneously
across all provinces and is absorbed by wave fixed effects.

% -------------------------------------------------------------
\subsection{Ordinal Measurement}\label{sec:ordinal_robust}
% -------------------------------------------------------------

The CFPS expectation question elicits direction (rise, same,
fall), not a point forecast.
The forecast-error coefficient $\hat{\beta}_3$ therefore
depends on the coding of the ordinal variable.
Three results confirm that the overreaction conclusion does
not depend on this coding choice.
The sign of $\hat{\beta}_3 < 0$ is invariant to any monotone
rescaling of the ordinal codes.
The pass-through coefficient $\hat{\beta}_2 = -5.25$, which
involves no ordinal variable, independently exceeds the
rational bound.
A Carlson-Parkin quantification that converts ordinal
responses to cardinal expected inflation rates confirms both
the sign and significance of the forecast-error coefficient.
An ordered probit that models the expectation response
without any cardinal assumption likewise confirms the positive
effect of meat shocks on expectations.
Appendix~\ref{app:scaling} provides full details.

% -------------------------------------------------------------
\subsection{Calibration Sensitivity}\label{sec:calib_robust}
% -------------------------------------------------------------

The rational benchmark depends on two parameters: meat's CPI
weight ($\omega$) and the persistence of the meat-price shock
($\rho$).
National estimates place $\omega$ in the range of 3 to 6
percent.
At the upper end ($\omega = 0.06$), the one-period rational
bound is 0.06; the time-aggregated bound
(Section~\ref{sec:time_aggregation}) is $24 \times 0.06 = 1.44$.
The estimated $\hat{\beta}_3 = -7.13$ exceeds even this
generous bound by a factor of five.

The persistence parameter $\hat{\rho} = 0.480$ is estimated
from a monthly AR(1) on province-level meat CPI with province
fixed effects (3,689 monthly observations).
The maximum rational coefficient obtains when the household's
persistence belief is maximally wrong, yielding
$|\beta_3^R| = \omega$.
No combination of $\omega$ and $\rho$ within economically
plausible ranges closes the gap between the rational bound
and the estimated coefficient.

